\expandafter\gdef\csname SNside1\endcsname{\node[anchor=north west,inner sep=0pt,text width=48mm,font=\SideFont,align=left] at ([xshift=138mm,yshift=-500.62pt]current page.north west) {\textcolor{SNink}{\textbf{课程连接 · 与组成原理按地址串联}}\par\vspace{1.5mm}组成原理研究Cache怎样定位已得到的存储地址，操作系统还要维护虚拟地址到物理页框的转换与权限。做综合访存题时先确认Cache使用虚拟还是物理字段，再排列TLB、页表、Cache和主存路径，避免混成一次查询。};}
\expandafter\gdef\csname SNside2\endcsname{\node[anchor=north west,inner sep=0pt,text width=48mm,font=\SideFont,align=left] at ([xshift=138mm,yshift=-397.86pt]current page.north west) {\textcolor{SNink}{\textbf{最小反例 · 空闲总量够仍不能分配}}\par\vspace{1.5mm}两个不相邻的空洞各为6MB，总空闲12MB；申请连续8MB仍失败，因为没有单个足够大的洞。只有物理相邻的空闲区才能直接合并，隔着占用区时需要搬移或改用非连续分配，不能把表中相邻行直接相加。};}
\expandafter\gdef\csname SNside3\endcsname{\node[anchor=north west,inner sep=0pt,text width=48mm,font=\SideFont,align=left] at ([xshift=138mm,yshift=-95.00pt]current page.north west) {\textcolor{SNink}{\textbf{不变量 · 映射替换页号不改偏移}}\par\vspace{1.5mm}页大小为\(2^d\)字节且页框同样大，转换前后低d位表示同一页内位置，保持不变；高位才由虚拟页号改成物理页框号。若计算结果的低d位发生变化，应先检查是否混入了页表项地址、字段位宽或容量单位。};}
\expandafter\gdef\csname SNside4\endcsname{\node[anchor=north west,inner sep=0pt,text width=48mm,font=\SideFont,align=left] at ([xshift=138mm,yshift=-95.00pt]current page.north west) {\textcolor{SNink}{\textbf{错项排除 · 多级表的收益来自稀疏}}\par\vspace{1.5mm}多级页表可不建立未使用虚拟区间对应的下级表，但每个有效映射仍要能表示页框和必要状态。若整个空间都使用，目录反而带来额外开销。因此不能说多级结构减少了每项位数，或对任何进程都一定节省总空间。};}
\expandafter\gdef\csname SNside5\endcsname{\node[anchor=north west,inner sep=0pt,text width=48mm,font=\SideFont,align=left] at ([xshift=138mm,yshift=-95.00pt]current page.north west) {\textcolor{SNink}{\textbf{课程连接 · 缓存对象不同导致分支不同}}\par\vspace{1.5mm}TLB缓存地址转换，数据Cache缓存内容。TLB没找到转换时，可从主存页表补入；数据仍可能已在Cache。缺页则要检查合法页面是否不驻留。把三种事件分别列条件，才能准确计算一次访问是否需要磁盘I/O。};}
\expandafter\gdef\csname SNside6\endcsname{\node[anchor=north west,inner sep=0pt,text width=48mm,font=\SideFont,align=left] at ([xshift=138mm,yshift=-95.00pt]current page.north west) {\textcolor{SNink}{\textbf{最小反例 · 驻留不解除越界检查}}\par\vspace{1.5mm}段2长300，已驻留且可写，访问偏移400仍越界，不能先算出基址加400再宣称访问成功。驻留位回答内容在哪里，权限回答能做什么，段长回答地址是否属于该段；三个检查条件互不替代。\par\smallskip\begin{tikzpicture}[x=1mm,y=1mm,>=Stealth,node distance=5mm]
\node[draw=SNink!50,fill=SNblue,rounded corners=1mm,text width=34mm,align=center,inner sep=1.2mm,font=\fontsize{7.5}{9}\selectfont] (m0) at (0,0) {段长300};
\node[draw=SNink!50,fill=SNblue,rounded corners=1mm,text width=34mm,align=center,inner sep=1.2mm,font=\fontsize{7.5}{9}\selectfont] (m1) at (0,-11) {偏移400};
\node[draw=SNink!50,fill=SNblue,rounded corners=1mm,text width=34mm,align=center,inner sep=1.2mm,font=\fontsize{7.5}{9}\selectfont] (m2) at (0,-22) {越界};
\draw[->,SNink] (m0) -- (m1);
\draw[->,SNink] (m1) -- (m2);
\end{tikzpicture}
\par{\fontsize{7.5}{9}\selectfont 段在内存且可写，也不能放行大于等于段长的偏移。}};}
\expandafter\gdef\csname SNside7\endcsname{\node[anchor=north west,inner sep=0pt,text width=48mm,font=\SideFont,align=left] at ([xshift=138mm,yshift=-95.00pt]current page.north west) {\textcolor{SNink}{\textbf{错项排除 · 修改位不是写权限}}\par\vspace{1.5mm}可写页面刚装入且只被读取时，可以写权限为1、修改位为0；只读页被写时，应产生保护异常，不能靠换入页面解决。修改位帮助判断淘汰是否回写，权限位限制行为，两项表示的状态不是同一件事。};}
\expandafter\gdef\csname SNside8\endcsname{\node[anchor=north west,inner sep=0pt,text width=48mm,font=\SideFont,align=left] at ([xshift=138mm,yshift=-520.73pt]current page.north west) {\textcolor{SNink}{\textbf{不变量 · 调页完成还要重新执行}}\par\vspace{1.5mm}合法页面不驻留时，原访存尚未成功。缺页处理把内容装入并更新映射后，通常回到原指令重试。因此费用表要区分失败查询、服务开销、重试访问；题目已包含的更新成本不能重复加，重试成本也不能漏掉。};}
\expandafter\gdef\csname SNside9\endcsname{\node[anchor=north west,inner sep=0pt,text width=48mm,font=\SideFont,align=left] at ([xshift=138mm,yshift=-228.38pt]current page.north west) {\textcolor{SNink}{\textbf{错项排除 · 固定分配约束页框数量}}\par\vspace{1.5mm}本节固定分配把各进程拥有的页框数保持为既定值；若允许直接淘汰别人的页并占为己用，就改变了双方数量。题目采用可变分配时才可讨论相应全局调整，不能把分配策略与置换范围的名字任意拼接。};}
\expandafter\gdef\csname SNside10\endcsname{\node[anchor=north west,inner sep=0pt,text width=48mm,font=\SideFont,align=left] at ([xshift=138mm,yshift=-95.00pt]current page.north west) {\textcolor{SNink}{\textbf{最小反例 · 命中也要更新LRU次序}}\par\vspace{1.5mm}两页框先访问A、B，再命中A，随后访问C。LRU应淘汰较久未用的B；FIFO仍可能淘汰最早装入的A。只在缺页时移动队列会遗漏命中对最近使用关系的影响，得到的就不是正确的LRU轨迹。};}
\expandafter\gdef\csname SNside11\endcsname{\node[anchor=north west,inner sep=0pt,text width=48mm,font=\SideFont,align=left] at ([xshift=138mm,yshift=-95.00pt]current page.north west) {\textcolor{SNink}{\textbf{课程连接 · 映射写入与磁盘持久化分开}}\par\vspace{1.5mm}程序通过映射地址修改内容，首先改变驻留页并可能置脏；何时写回文件还取决于映射属性和写回机制。文件系统维护后备关系，虚拟内存维护映射和驻留，二者协作不代表每一条赋值都会立即产生一次磁盘写入。};}
\expandafter\gdef\csname SNside12\endcsname{\node[anchor=north west,inner sep=0pt,text width=48mm,font=\SideFont,align=left] at ([xshift=138mm,yshift=-95.00pt]current page.north west) {\textcolor{SNink}{\textbf{不变量 · 窗口先取引用再去重}}\par\vspace{1.5mm}工作集窗口若规定最近6次引用，应先截出准确的6次，再对页号去重。同一页出现三次仍只占集合中一项。不能先找6个不同页，也不能从当前时刻向未来补引用，否则集合不再代表题设的近期需求。};}
\expandafter\gdef\csname SNside13\endcsname{\node[anchor=north west,inner sep=0pt,text width=48mm,font=\SideFont,align=left] at ([xshift=138mm,yshift=-95.00pt]current page.north west) {\textcolor{SNink}{\textbf{不变量 · 窗口先取引用再去重}}\par\vspace{1.5mm}工作集窗口若规定最近6次引用，应先截出准确的6次，再对页号去重。同一页出现三次仍只占集合中一项。不能先找6个不同页，也不能从当前时刻向未来补引用，否则集合不再代表题设的近期需求。};}
\expandafter\gdef\csname SNside14\endcsname{\node[anchor=north west,inner sep=0pt,text width=48mm,font=\SideFont,align=left] at ([xshift=138mm,yshift=-215.38pt]current page.north west) {\textcolor{SNink}{\textbf{课程连接 · 映射写入与磁盘持久化分开}}\par\vspace{1.5mm}程序通过映射地址修改内容，首先改变驻留页并可能置脏；何时写回文件还取决于映射属性和写回机制。文件系统维护后备关系，虚拟内存维护映射和驻留，二者协作不代表每一条赋值都会立即产生一次磁盘写入。};}
\expandafter\gdef\csname SNside15\endcsname{\node[anchor=north west,inner sep=0pt,text width=48mm,font=\SideFont,align=left] at ([xshift=138mm,yshift=-180.56pt]current page.north west) {\textcolor{SNink}{\textbf{最小反例 · 命中也要更新LRU次序}}\par\vspace{1.5mm}两页框先访问A、B，再命中A，随后访问C。LRU应淘汰较久未用的B；FIFO仍可能淘汰最早装入的A。只在缺页时移动队列会遗漏命中对最近使用关系的影响，得到的就不是正确的LRU轨迹。};}
\expandafter\gdef\csname SNside16\endcsname{\node[anchor=north west,inner sep=0pt,text width=48mm,font=\SideFont,align=left] at ([xshift=138mm,yshift=-95.00pt]current page.north west) {\textcolor{SNink}{\textbf{错项排除 · 固定分配约束页框数量}}\par\vspace{1.5mm}本节固定分配把各进程拥有的页框数保持为既定值；若允许直接淘汰别人的页并占为己用，就改变了双方数量。题目采用可变分配时才可讨论相应全局调整，不能把分配策略与置换范围的名字任意拼接。};}
\AddToShipoutPictureFG{\begin{tikzpicture}[remember picture,overlay]\ifcsname SNside\thepage\endcsname\csname SNside\thepage\endcsname\fi\end{tikzpicture}}